\section{Discussion and Limitations}
\label{sec:discussion}

\subsection{What a support desk should take from this}

\textbf{Consume the posterior, not the argmax.} The single largest effect we
measure comes from changing how an existing model's output is read, at no
training cost: the argmax-tier policy creates a seven-minute cross-script
disparity that the same model's posterior largely removes.

\textbf{Audit triage fairness in waiting time, not in F\textsubscript{1}.} A
per-group accuracy table would have shown a few points of deficit on romanized
tracks and given no sense of what it costs a customer. The same deficit,
denominated in minutes and controlled against FIFO and a gold-label oracle,
names both a cost and a cause.

\textbf{Better classification is not the next thing to buy.} Our score attains
99\% of what perfect labels would deliver, and substituting a weaker classifier
leaves the median untouched. Effort is better spent on how predictions are
consumed, and on the tail behaviour that F\textsubscript{1} hides.

\subsection{Limitations}

\textbf{The simulation parameters are assumed.} Arrival rate, service
distribution, agent count, and the service windows are stated assumptions rather
than measurements from a live desk. This is the largest threat to the paper. We
mitigate it by sweeping utilisation and by reporting attainment, which is a ratio
and so is less sensitive to the absolute scale than raw waiting times, but a
parameterisation drawn from published helpdesk statistics would be stronger.
\gap{Arrival and service statistics from a real or published helpdesk.}

\textbf{The romanized tracks are synthetic}, generated from the native tracks
rather than typed by people \cite{chakravarthi2020tamilmix}. Real romanized input
is harder, so the disparity we measure is a conservative estimate.

\textbf{The Tanglish rendering is defective}, at 60.3\% out-of-vocabulary against
its own training split. Singlish carries the disparity result independently.

\textbf{Two label columns are model-generated}, with ceilings of 0.7931 and
0.7722 reported throughout, and the sentiment ceiling rests on one annotator and
31 negative tickets.

\textbf{Confidence intervals are within-fit.} Classification intervals resample
the test set rather than the training run, and one seed repeat moved a sentiment
headline by 0.0158. Queue intervals resample simulation seeds and are not subject
to this.

\textbf{The weight vector is not identified.} The development surface is flat
across $w_P \in [0.10, 0.90]$; we claim only that priority dominates and that no
corner is competitive.
